\subsection{Settings}

\subsubsection{Metrics}

To quantify the performance of the algorithm, we use different metrics tailored
to each query task.
Let $\hat\theta$ denote an estimator of the true value $\theta$, and let
$\hat\theta_i$ denote the estimate of the true value $\theta_i$ in the $i$-th query,
for $i=1,\ldots,n_q$.
\begin{itemize}
    \item {Mean Absolute Error (MAE)}.
MAE quantifies the average magnitude of estimation errors and serves as a
standard indicator of overall accuracy across all query types.
The MAE of an estimator $\hat\theta$ is defined as
\[
  \mathsf{MAE}(\hat\theta)
  \triangleq \E(|\hat\theta - \theta |)
  = \frac{1}{n_q}\sum_{i=1}^{n_q} |\hat\theta_i - \theta_i |.
\]
    \item {Mean Relative Error (MRE)}.
MRE measures the error relative to the true value, offering
a scale-invariant perspective on accuracy.
The MRE of an estimator $\hat\theta$ is defined as
\[
  \mathsf{MRE}(\hat\theta)
  \triangleq \E(\frac{|\hat\theta - \theta |}{\theta})
  = \frac{1}{n_q}\sum_{i=1}^{n_q}\frac{|\hat\theta_i - \theta_i|}{\theta_i}.
\]
    \item {Receiver Operating Characteristic (ROC) Curve}.
For event monitoring, we use the ROC curve to evaluate the trade-off between
true positive and false positive rates across varying decision thresholds.
The Area Under the Curve (AUC) is used as a key metric for evaluating the
model's overall performance.
A higher AUC indicates better event detection performance.
\end{itemize}
% \bullethdr{Mean Absolute Error (MAE)}.
% MAE quantifies the average magnitude of estimation errors and serves as a
% standard indicator of overall accuracy across all query types.
% The MAE of an estimator $\hat\theta$ is defined as
% \[
%   \mathsf{MAE}(\hat\theta)
%   \triangleq \E(|\hat\theta - \theta |)
%   = \frac{1}{n_q}\sum_{i=1}^{n_q} |\hat\theta_i - \theta_i |.
% \]

% \bullethdr{Mean Relative Error (MRE)}.
% MRE measures the error relative to the true value, offering
% a scale-invariant perspective on accuracy.
% The MRE of an estimator $\hat\theta$ is defined as
% \[
%   \mathsf{MRE}(\hat\theta)
%   \triangleq \E(\frac{|\hat\theta - \theta |}{\theta})
%   = \frac{1}{n_q}\sum_{i=1}^{n_q}\frac{|\hat\theta_i - \theta_i|}{\theta_i}.
% \]

% \bullethdr{Receiver Operating Characteristic (ROC) Curve}.
% For event monitoring, we use the ROC curve to evaluate the trade-off between
% true positive and false positive rates across varying decision thresholds.
% The Area Under the Curve (AUC) is used as a key metric for evaluating the
% model's overall performance.
% A higher AUC indicates better event detection performance.

\subsubsection{Baselines}

We compare MTSP-LDP with several state-of-the-art $w$-event LDP algorithms,
including LBA, LBD, LSP, and LBU, as introduced in \cref{ss:competitors}.
To ensure fair comparison and reproducibility, all methods, including MTSP-LDP,
were implemented in Python under a unified experimental framework.
All experiments were conducted on a desktop equipped with an Intel Core i7-10700
CPU (2.90 GHz) and 16 GB of RAM.
